% Original: PDF strana 20, gornji slajd.
\slajdid{02-s039}
\begin{frame}[label=02-s039]{Površine reda \(k\) u Karnoovoj karti}
% element: 02-s039:e01
\begin{center}
\begingroup\setlength{\tabcolsep}{10pt}
\begin{tabular}{c|rrrrrrr}
\toprule
Red \(k\) & 0 & 1 & 2 & 3 & 4 & 5 & 6 \\
\midrule
Broj susednih polja \(2^k\) & 1 & 2 & 4 & 8 & 16 & 32 & 64 \\
\bottomrule
\end{tabular}
\endgroup
\end{center}
\par\vspace{6mm}
% element: 02-s039:e02
\par\Rezultat{Površina mora biti pravilna: pravougaonik ili kvadrat u ravni, kvadar ili kocka u 3D, odgovarajući hiperkvadar u višoj dimenziji.}
\end{frame}
